%%%%%%%%%%%%%%%%%%%%%%%%%%%%%%%%%%%%%%%%%%%%%%%%%%%%%%%%%%%%%%%%%%%%%%%%%%%%%%%%%%%%%%%%%%%%%%%%%%%%%%
%
%   Filename    : abstract.tex 
%
%   Description : This file will contain your abstract.
%                 
%%%%%%%%%%%%%%%%%%%%%%%%%%%%%%%%%%%%%%%%%%%%%%%%%%%%%%%%%%%%%%%%%%%%%%%%%%%%%%%%%%%%%%%%%%%%%%%%%%%%%%

\begin{abstract}
Social communication can be a challenge for many a youth, especially those with neurodivergent traits or social anxiety. This can ostensibly be linked to difficulty with pragmatic competence, which often goes beyond the most shallow and direct meanings of what is said in conversation. Despite many recent improvements in NLP that help AI chatbots seem to converse naturally, they still face difficulty with pragmatic communication. This paper proposes the development of Socius, a chatbot aimed to hold pragmatically-aware conversations to enhance one-on-one conversation skills in youth. The system operates within the APAM (AI Partner - AI Mentor) framework for the delivery of support informed by Grice’s Maxims. Socius is to help users practice conversation in realistic, simulated scenarios while also reflecting on relevant pragmatic theories such as relevance, clarity, and implicature. The system is build with Python and OpenAI’s GPT-4 API, trained on the GRICE dataset and augmented with synthetic examples. The study will involve knowledge base construction, software development, user data collection, and qualitative validation. By offering a robust and safe grounded-theory tool in Socius, this paper seeks to offer a novel approach in personalized communication training and broader efforts in inclusive conversational AI design.%
%  Do not put citations or quotes in the abract.
%



\begin{flushleft}
\begin{tabular}{lp{4.25in}}
\hspace{-0.5em}\textbf{Keywords:}\hspace{0.25em} & Human-centered computing, Human computer interaction (HCI), HCI design and evaluation methods, Natural language processing, Conversational AI, Assistive technologies,  Education technologies
\end{tabular}
\end{flushleft}
\end{abstract}
